%=====================================================================
\chapter{Publishing: Venue, Submission and Peer Review}\label{ch:publishing}
%=====================================================================
\locator{6}{Part~6 --- Producing and Defending Research}

\begin{outcomes}
By the end of this chapter you will be able to:
\begin{tight}
  \item \textbf{Select} a venue on fit as well as on category.
  \item \textbf{Write} a cover letter that does its one job.
  \item \textbf{Respond} to reviewers so that a revision is accepted.
  \item \textbf{Handle} rejection as a step rather than a verdict.
  \item \textbf{Plan} a publication sequence that satisfies \S15 without
    slicing one study into two.
\end{tight}
\end{outcomes}

\begin{prereq}
\chref{publication} for the venue landscape and the HJRS categories;
\chref{writing} for the manuscript; \chref{integrity} for authorship and text
recycling.
\end{prereq}

\begin{keyterms}
cover letter $\bullet$ desk rejection $\bullet$ major and minor revision
$\bullet$ response letter $\bullet$ rebuttal $\bullet$ appeal $\bullet$
reviewer 2 $\bullet$ resubmission $\bullet$ publication sequence
\end{keyterms}

\section{Choosing Where to Send It}

\begin{center}
\small
\begin{tabular}{@{}L{3.4cm}L{10.9cm}@{}}
\toprule
Category floor   & HJRS W, X or Y as \S15 requires. Verify on the day, and keep
                   a dated record (\chref{publication}) \\
Fit              & Does the venue publish work you cite? This predicts
                   acceptance and readership better than any other signal \\
Scope            & Read the aims and scope, and three recent papers. Most desk
                   rejections are scope mismatches \\
Length and format& A 12-page study does not fit a 6-page venue without damage \\
Review model     & Double-blind requires anonymising; plan for it before
                   writing \\
Time to decision & Ask colleagues, or check the journal's reported medians.
                   This matters when \S15 sits between you and submission \\
Cost             & Any APC, and whether you can pay it (\chref{publication}) \\
\bottomrule
\end{tabular}
\end{center}

\begin{conceptbox}[Aim one level above where you expect to land]
Review at a stronger venue is usually more thorough, and a rejection with three
substantive reviews improves the paper before it goes elsewhere. The cost is
time, and the calculation changes when a deadline binds --- but for a first
paper with time in hand, one attempt at the better venue is usually worth it.

\medskip
What is not worth it is submitting to a venue whose scope you do not match, in
hope. That is a desk rejection, which costs weeks and returns no reviews at
all.
\end{conceptbox}

\section{The Cover Letter}

Its job is to tell the editor why this paper belongs in this venue. Nothing
else.

\begin{center}
\small
\begin{tabular}{@{}L{0.6cm}L{13.8cm}@{}}
\toprule
1 & Title, type of submission, and the venue's own scope area it addresses \\
2 & The contribution, in two sentences \\
3 & Why this venue --- ideally naming papers it has published that this
    extends \\
4 & Declarations: original work, not under review elsewhere, any prior
    conference version (\chref{integrity}), any conflicts, ethics approval \\
5 & Suggested or excluded reviewers, if the venue asks \\
\bottomrule
\end{tabular}
\end{center}

\begin{pitfall}[What not to put in a cover letter]
Do not summarise the abstract --- the editor has it. Do not claim the work is
groundbreaking; editors read that in most letters and discount it. Do not
explain why the topic is important at length. Keep it under one page, and let
the paper do the work.
\end{pitfall}

\section{Responding to Reviewers}

This is a distinct skill and it is worth more than most graduate scholars
realise: a good response letter converts a major revision into an acceptance,
and a bad one converts it into a rejection.

\begin{center}
\small
\begin{tabular}{@{}L{3.4cm}L{10.9cm}@{}}
\toprule
Answer everything  & Every comment gets a response, including the ones you
                     disagree with. An unanswered comment reads as evasion \\
Quote the comment  & Then respond beneath it. The editor is checking coverage
                     and should not have to search \\
Say what you changed & And where: ``revised, Section 4.2, lines 310--318''.
                     Point to the text, do not describe it \\
Disagree politely and with evidence & ``We considered this and retained the
                     original approach because \ldots'' followed by a reason,
                     not an assertion \\
Thank them once    & At the top. Not after every point \\
Highlight changes  & A marked-up version alongside the clean one \\
\bottomrule
\end{tabular}
\end{center}

\begin{templatebox}[C16 --- Reviewer response letter]
\textbf{Opening.} Thank the editor and reviewers; state the main changes in
three or four sentences so the editor can see the shape of the revision without
reading the whole letter.
\filllines{3}

\medskip
\noindent Then, for every single comment, repeat this block:

\medskip
\noindent\textbf{Comment R1.1} --- quoted verbatim:
\filllines{1}

\noindent\textbf{Response:}
\filllines{2}

\noindent\textbf{Change made} --- section and line numbers, not a description:
\filllines{1}

\medskip
\noindent Where you disagree, use this form: \emph{``We understand the concern
that [restate it accurately]. We have retained [X] because [reason]. To address
the underlying issue we have [what you did instead].''} Conceding something is
almost always available, and it makes the disagreement land better.
\end{templatebox}

\begin{alertbox}[Restate the objection before answering it]
The single most effective move in a response letter: begin each disputed point
by stating the reviewer's concern in your own words, accurately and
sympathetically, before disagreeing.

\medskip
It proves you understood rather than dismissed them, and reviewers who feel
understood are markedly more willing to be persuaded. It is the same discipline
as charitable reading (\chref{appraisal}), applied from the other side of the
process.
\end{alertbox}

\begin{conceptbox}[When a reviewer is simply wrong]
It happens. They misread a section, or they are not expert in the specific
method.

\medskip
Even then, the misreading is data: if a competent reviewer misread it, other
readers will too, so something in the writing invited it. Respond by clarifying
the text as well as correcting the reviewer, and say that is what you have
done. This turns a dispute into an improvement, and editors --- who must decide
between you --- notice which party is doing that.
\end{conceptbox}

\section{Outcomes}

\begin{center}
\small
\begin{tabular}{@{}L{3.4cm}L{10.9cm}@{}}
\toprule
Accept              & Rare on first submission. Check the proofs carefully;
                      errors introduced at typesetting are yours to catch \\
Minor revision      & Effectively acceptance. Do the work quickly and
                      thoroughly \\
Major revision      & A real opportunity: most major revisions that respond
                      seriously are eventually accepted. Do not treat it as a
                      soft rejection \\
Reject and resubmit & A new submission with the review history attached.
                      Substantial work, but the door is open \\
Reject              & Read the reviews, wait a few days, then improve the paper
                      and send it elsewhere \\
Desk reject         & Usually scope or format. Fix and send elsewhere quickly;
                      it costs little \\
\bottomrule
\end{tabular}
\end{center}

\begin{conceptbox}[Rejection is the normal case]
Selective venues reject three-quarters or more of submissions. Most published
papers were rejected somewhere first, by authors far more experienced than you.

\medskip
The professional response is mechanical, and worth making a habit: read the
reviews once, put them aside for two or three days, then extract every
actionable criticism into a list --- including from the review that felt
unfair. Fix what is fixable, note what is not, and submit elsewhere within two
weeks. Papers that sit after a rejection tend not to be resubmitted at all,
and that is how work is lost.
\end{conceptbox}

\section{Planning the Sequence}

\begin{regbox}[The constraint you are planning around --- PhD \S15]
One W-category paper, or two X-category papers, for Science disciplines. First
author. Relevant to the dissertation. \textbf{Published after synopsis
approval.} MS scholars have no publication requirement.
\end{regbox}

\begin{center}
\small
\begin{tabular}{@{}L{3.0cm}L{5.4cm}L{5.5cm}@{}}
\toprule
\textbf{Stage} & \textbf{Output} & \textbf{Timing} \\
\midrule
After the review
  & The systematic review, as a standalone paper (\chref{slr})
  & Submit after synopsis approval. Reviews are well suited to journals, and
    this is a genuine second contribution rather than a slice \\
After study 1
  & The empirical or design-science study
  & The main paper. Conference first if the fit is right, then a substantially
    extended journal version \\
Throughout
  & The artefact, deposited with a DOI (\chref{reproducibility})
  & Citable in its own right, and it strengthens both papers \\
\bottomrule
\end{tabular}
\end{center}

\begin{alertbox}[The sequencing trap, stated once more]
A paper accepted \emph{before} your synopsis is approved does not count towards
\S15. Scholars arriving with work already in progress lose it this way
routinely.

\medskip
If you have unpublished work in hand, get the synopsis approved first, then
submit. The delay is a few months; the alternative is a paper that does not
count and cannot be republished.
\end{alertbox}

\begin{traditions}{publishing}
\tradEMP{Journals and archival conferences. Reviewers increasingly expect data
and analysis scripts alongside the manuscript.}
\tradDSR{Information systems and design-oriented venues; the artefact is
submitted for badging alongside the paper (\chref{reproducibility}).}
\tradQUAL{Venues that allow the length the method requires. A page limit that
forces the removal of thick description damages the work rather than compressing
it.}
\tradFORM{Theory venues with full proofs in appendices; correctness is
refereed in detail and the review cycle is correspondingly slow.}
\end{traditions}

\begin{lab}[Prepare a submission]
\begin{tightnum}
  \item Shortlist three venues. For each record HJRS category with the date
    checked, scope fit, format, review model, typical decision time and cost.
  \item Count how many papers you cite that appeared in each. Rank by that
    number.
  \item Write the one-page cover letter for your first choice.
  \item Take any set of review comments --- your supervisor's on a draft will
    do --- and write a full response letter using template C16, including one
    point where you disagree.
  \item Draw your publication sequence against your degree timeline, marking
    synopsis approval and showing that every counted paper falls after it.
\end{tightnum}
\end{lab}

\begin{checkpoint}
\begin{tightnum}
  \item Why is the count of your own citations appearing in a venue a better
    signal than its impact factor?
  \item What is the one job of a cover letter, and what should be left out?
  \item A reviewer misreads your method. What do you do, and why does
    correcting only the reviewer waste the opportunity?
  \item You have a paper accepted two months before your synopsis is approved.
    What does \S15 say, and what should you have done?
\end{tightnum}
\end{checkpoint}

\begin{chaptersummary}
\begin{tight}
  \item Clear the HJRS category floor, then choose on fit. The venues
    publishing work you cite are where your readers are.
  \item The cover letter says why this paper belongs in this venue. One page.
  \item Answer every comment, quote it, point to the changed lines, and restate
    an objection before disagreeing with it.
  \item Major revision is an opportunity, not a soft rejection. Rejection is
    the normal case; extract the actionable points and resubmit within two
    weeks.
  \item Plan the sequence: review paper, then study paper, both after synopsis
    approval. Two genuine contributions, not one sliced in half.
\end{tight}
\end{chaptersummary}

\begin{reviewq}
  \item Response letters are rarely taught and heavily weighted in editorial
    decisions. Why is this skill invisible in graduate training, and what would
    it take to teach it?
  \item Aiming above your expected level buys better reviews at the cost of
    time. For a PhD scholar under \S15's sequencing constraint, model that
    trade-off concretely.
  \item Reviewer anonymity permits honesty and enables carelessness. Would open
    review improve computing research? Consider the effect on junior reviewers
    specifically.
  \item \S15 requires publication before dissertation submission, which places
    a scholar's graduation at the mercy of journal timelines they do not
    control. Propose an alternative that preserves the requirement's intent.
\end{reviewq}
